\clearpage
\begin{proofbox}
\textbf{\textcolor{CProof}{先发现规律：从具体函数一步一步走到抽象函数}}

取 $f(t)=\sqrt{t-1}+\sqrt{3-t}$。先不谈新规则，
只做初中就学过的“根号里的数不能为负”。
\[
 t-1\ge0,\quad3-t\ge0
 \quad\Longleftrightarrow\quad1\le t\le3.
\]
把 $t$ 换成 $2x-1$，仍然满足原来的两条要求：
\[
\left\{\begin{aligned}
 &(2x-1)-1\ge0,\\
 &3-(2x-1)\ge0
\end{aligned}\right.
\quad\Longleftrightarrow\quad 1\le x\le2.
\]
把上面的两条根号条件合并来看，恰好就是
\[
 1\le \underbrace{2x-1}_{\text{实际送进 }f\text{ 的数}}\le3.
\]
\textbf{观察：}我们没有改变 $f$ 的计算方法，
只改变了送进计算方法的数。
现在即便不告诉你 $f$ 的式子，只告诉你 $f$ 能接收 $[1,3]$，
同样可以直接解 $1\le2x-1\le3$。
从具体计算归纳出来的做法，正是下一节方法一、方法二的依据。

\medskip

\textbf{\textcolor{CProof}{方法一：像检查入场券一样推出来（最容易理解）}}
\begin{center}\begin{tikzpicture}[font=\small,>=stealth,line join=round]
\draw[gray!45,rounded corners=3pt] (0,0) rectangle (14.2,9.4);
\node at (7.1,8.8) {把允许的输入范围拉回到自变量轴};

\begin{scope}[shift={(2.2,2.6)},scale=1.35]
\fill[blue!8] (-.3,1) rectangle (2.6,3);
\draw[->] (-.5,0)--(2.8,0) node[right]{$x$};\draw[->] (0,-.7)--(0,3.6) node[above]{$t$};
\draw[teal!80!black,thick,domain=.2:2.2] plot (\x,{2*\x-1}) node[above right]{$t=2x-1$};
\foreach \x/\y in {1/1,2/3}{\draw[dashed,gray] (\x,0)--(\x,\y)--(0,\y);\fill (\x,\y) circle(1.5pt);\node[below] at (\x,0){$\x$};\node[left] at (0,\y){$\y$};}
\draw[blue,line width=2pt] (0,1)--(0,3);\draw[orange!90!black,line width=2pt] (1,0)--(2,0);
\end{scope}
\node[align=left,anchor=west] at (8,5.1) {外层允许：$t\in[1,3]$\\[8pt] 内层关系：$t=2x-1$\\[8pt] $1\le2x-1\le3$\\[8pt] 自变量：$x\in[1,2]$};
\node[text width=13.2cm,align=center] at (7.1,.65) {先检查括号里的数，再解出 $x$ 的范围。};
\end{tikzpicture}
\end{center}


设 $f$ 只允许接收 $[1,3]$ 内的数。现在要研究 $f(2x-1)$。

\begin{enumerate}[leftmargin=*,itemsep=0.42em]
 \item \textbf{先看最后一步。} 只有送进 $f$ 的数在 $[1,3]$ 内，函数值才有意义。
 \item \textbf{再看中间一步。} 送进 $f$ 的是 $2x-1$，所以它必须满足
 \[
 1\le2x-1\le3.
 \]
 \item \textbf{找到最开始的数。} 三个部分同时加 $1$，再除以 $2$：
 \[
 2\le2x\le4\quad\Longleftrightarrow\quad1\le x\le2.
 \]
 因此 $f(2x-1)$ 的定义域是 $[1,2]$。
 \item \textbf{倒过来检查。} 如果已知 $f(2x-1)$ 的完整定义域就是 $[1,2]$，
 那么 $2x-1$ 从 $1$ 变化到 $3$，中间的数也全部取到。
 而给定任何实数 $t$，总能取 $x=(t+1)/2$ 使 $2x-1=t$。
 所以没有“永远无法送进 $f$ 的实数”，便可断定 $f(t)$ 的完整定义域是 $[1,3]$。
\end{enumerate}

\textbf{还能直接看懂反过来的道理：}
如果已经知道 $f(2x-1)$ 的\textbf{完整}定义域是 $[1,2]$，
那么它实际送出的数是 $[1,3]$。
对任何实数 $t$，都能用 $x=(t+1)/2$ 造出 $2x-1=t$。
这样就没有遗漏的输入数，因而可以把 $[1,3]$ 当作 $f$ 的完整定义域。
如果括号里是 $x^2$，负数永远造不出来，这一步便不能照搬。

\medskip

\textbf{为什么成立？}\quad 一个数是否可以使用，不取决于 $f$ 长什么样，
只取决于真正送进去的数是否被允许。

\medskip
\end{proofbox}
\clearpage
\begin{proofbox}
\textbf{\textcolor{CProof}{方法二：换一个字母，给出通用推导（常规方法）}}
\begin{center}\begin{tikzpicture}[font=\small,>=stealth,line join=round]
\draw[gray!45,rounded corners=3pt] (0,0) rectangle (14.2,7.2);
\node at (7.1,6.6) {反求完整定义域：每个实数输入都有唯一原像};

\draw[->] (1.4,4.9)--(12.7,4.9) node[right]{$x$};\draw[->] (1.4,2.6)--(12.7,2.6) node[right]{$t$};
\draw[orange!90!black,line width=2pt] (3,4.9)--(7.4,4.9);\draw[blue,line width=2pt] (3,2.6)--(7.4,2.6);
\foreach \u/\v/\pos in {1/1/3,2/3/7.4}{\fill (\pos,4.9) circle(2pt);\fill (\pos,2.6) circle(2pt);\node[above] at (\pos,4.9){$\u$};\node[below] at (\pos,2.6){$\v$};\draw[->,teal!80!black,thick] (\pos,4.6)--(\pos,2.9);}
\node at (5.2,5.8){完整定义域 $A=[1,2]$};\node at (5.2,1.8){$D_f=[1,3]$};
\node[align=left] at (10.2,3.8){正向 $t=2x-1$\\[10pt] 反向 $x=(t+1)/2$};
\node[text width=13.2cm,align=center] at (7.1,.65) {一般地，$a\ne0$ 时 $t=ax+b$ 与 $x=(t-b)/a$ 互相对应。};
\end{tikzpicture}
\end{center}


\textbf{第一部分：已知 $f(t)$ 的定义域，求 $f[g(x)]$。}

设 $f(t)$ 能接收的全部数构成 $D$，把 $g(x)$ 记成 $t$。
当 $g(x)$ 有意义时，
\[
 f[g(x)]\text{有意义}
 \quad\Longleftrightarrow\quad
 f(t)\text{有意义}
 \quad\Longleftrightarrow\quad t\in D.
\]
把 $t=g(x)$ 放回去，就得到
\[
 \boxed{\ g(x)\text{有意义，且 }g(x)\in D\ }.
\]
最后解出 $x$ 的范围，就是所求定义域。

\textbf{第二部分：已知 $f(ax+b)$ 的完整定义域 $A$，反求 $f(t)$，其中 $a\ne0$。}

要说明 $f$ 能接收的数\textbf{恰好}是 $\{ax+b\mid x\in A\}$，需要检查两个方向。
\begin{enumerate}[leftmargin=*,itemsep=0.35em]
 \item 只要 $x\in A$，题目就保证 $f(ax+b)$ 有意义，
 所以 $t=ax+b$ 一定是 $f$ 能接收的数。
 \item 反过来，任取一个 $f$ 能接收的数 $t$。
 因为 $a\ne0$，可以解出唯一的
 \[
 x=\frac{t-b}{a},\qquad ax+b=t.
 \]
 此时 $f(ax+b)=f(t)$ 有意义，因此这个 $x$ 必须在题目给出的\textbf{完整定义域} $A$ 中。
 于是 $t=ax+b$ 一定来自某个 $x\in A$。
\end{enumerate}
两个方向合起来，没有多，也没有少：
\[
 \boxed{\ f(t)\text{的定义域}=\{ax+b\mid x\in A\}\quad(a\ne0).\ }
\]
\textbf{特别注意：} $a\ne0$ 非常重要；若 $a=0$，$ax+b$ 永远只有 $b$ 一个数，
我们不能由 $f(b)$ 是否有意义推断 $f$ 对其他数如何。

\medskip
\end{proofbox}
\clearpage
\begin{proofbox}
\textbf{\textcolor{CProof}{方法三：用两个具体函数检验反推的限制（辨别误用）}}
\begin{center}\begin{tikzpicture}[font=\small,>=stealth,line join=round]
\draw[gray!45,rounded corners=3pt] (0,0) rectangle (14.2,9.4);
\node at (7.1,8.8) {平方只能探测非负输入，负半轴仍然未知};

\begin{scope}[shift={(3.7,2.6)},x=1cm,y=.95cm]
\fill[gray!12] (-2.6,-.9) rectangle(2.6,0);
\draw[->] (-2.7,0)--(2.8,0) node[right]{$x$};\draw[->] (0,-.2)--(0,5.1) node[above]{$t$};
\draw[teal!80!black,thick,domain=-2.15:2.15,samples=80] plot(\x,{\x*\x});
\draw[dashed] (-2,0)--(-2,4)--(2,4)--(2,0);\draw[orange!90!black,line width=2pt] (-2,0)--(2,0);\draw[blue,line width=2pt] (0,0)--(0,4);
\node[below] at (-2,0){$-2$};\node[below] at (2,0){$2$};\node[above left] at (0,4){$4$};\node[anchor=west] at (1.3,.9){$t=x^2$};\node at (0,-.6){$t<0$ 无法取到};
\end{scope}
\node[align=left,anchor=west] at (7.9,5.1){同一复合定义域 $[-2,2]$：\\[10pt] $D_{f_1}=[0,4]$\\[8pt] $D_{f_2}=(-\infty,4]$\\[12pt] 共同的可探测部分\\[4pt] $D_f\cap[0,+\infty)=[0,4]$};
\node[text width=13.2cm,align=center] at (7.1,.65) {两个外层函数在负数上不同，复合后却无法区分。};
\end{tikzpicture}
\end{center}


假设有人说：“已知 $f(x^2)$ 的定义域是 $[-2,2]$，
所以 $f(t)$ 的定义域必为 $[0,4]$。”
我们来找两个完全符合题目、却有不同定义域的函数。

\textbf{第一种：}
\[
 f_1(t)=\sqrt{t(4-t)}.
\]
根号里的数必须不小于 $0$，所以 $f_1(t)$ 的定义域是 $[0,4]$。
代入 $t=x^2$，有意义的条件是 $0\le x^2\le4$，即 $-2\le x\le2$。

\textbf{第二种：}
\[
 f_2(t)=\sqrt{4-t}.
\]
只需 $4-t\ge0$，所以 $f_2(t)$ 的定义域是 $(-\infty,4]$。
代入 $t=x^2$，有意义的条件还是 $x^2\le4$，
即 $-2\le x\le2$。

两种函数的 $f(x^2)$ 定义域\textbf{完全相同}，
但 $f(t)$ 的定义域\textbf{不同}！
原因很简单：$x^2$ 不可能是负数，负数有没有被 $f$ 接受，从来没有被试出来。

\medskip
\textbf{\textcolor{CProof}{最终判断}}

正着求：检查送进去的数是否被允许，始终成立。
反着求完整定义域：必须确保括号里的式子能够取遍所有实数；
否则通常只能知道 $f$ 在\textbf{实际送得到的那些数}上有没有意义。
\end{proofbox}
